% Lösungen Einheit 4 von 5 – Gleichung bestimmen (zusammenbau v0.1)
\einheitenkopf{Einheit 4 von 5 · Gleichung bestimmen}

\erg{24}{a) $m = 2$, $n = -3$: $f(x) = 2x - 3$ \quad b) $7 = 2 \cdot 4 + n$, $n = -1$: $f(x) = 2x - 1$ \quad c) $m = (10 - 4) : (2 - 0) = 3$, $n = 4$: $f(x) = 3x + 4$ \quad d) $m = (9 - 3) : (5 - 2) = 2$, $n = -1$: $f(x) = 2x - 1$ \quad e) $m = (0 - 6) : (3 - 1) = -3$, $n = 9$: $f(x) = -3x + 9$ \quad f) Gerade durch A und B; erste Aussage falsch, zweite wahr; $m = -1{,}5$, $n = 5$: $f(x) = -1{,}5x + 5$}
\erg{25}{a) Gerade durch $A(-3|-1)$ und $B(2|4)$}
\erg{26}{a) $3x - 4 = -x + 8$, $x = 3$, $y = 5$: $S(3|5)$}
\erg{27}{a) Kim rechnet oben A minus B, unten B minus A; richtig: $m = (12 - 3) : (5 - 2) = 3$. \quad b) Nein: Beide haben dieselbe Steigung, aber verschiedene y-Achsenabschnitte; sie sind parallel. \quad c) $m = -0{,}5$, $n = 1$: $f(x) = -0{,}5x + 1$ \quad d) $m = (25{,}20 - 12{,}60) : (10 - 4) = 2{,}10$; $n = 12{,}60 - 4 \cdot 2{,}10 = 4{,}20$: Grundgebühr 4{,}20 €}
